% Source-preserving integrated proof body; input from toe_round7_proofs.tex.
\subsubsection*{Round 7: a microscopic flux-string disorder operator for the frozen QWZ seam}
\label{sec:toe-r7-charged}
\noindent\emph{The displayed arguments below are the proof-bearing source.  The current
repository regression is \veri{v1033\_\allowbreak charged\_\allowbreak disorder.py}; its finite checks audit the
implementation and do not replace the universal steps in the proof.}

\begingroup
\def\Hh{\mathcal H}
\def\F{\mathcal F}
\def\C{\mathbb C}
\def\Z{\mathbb Z}
\def\id{\mathbf 1}
\def\norm#1{\left\lVert #1\right\rVert}
\begin{keybox}[Scope and exact claim boundary]
There is a concrete microscopic boundary-condition-changing operator for the
$M=1$, $N_y=8$ QWZ seam once one local four-level flux register is adjoined.
Couple its clock $Z$ directly to the seam hopping and let its shift $X$ raise
the flux.  Dressing $X$ by a quarter charge rotation on a fermion arc gives a
disorder string $D_A$.  It has nonzero $r\to r+1$ corners, order four, commutes
with canonical fermion parity, and its commutator with the actual QWZ
Hamiltonian is supported exactly on the opposite endpoint bond.  Thus this is
not the earlier arbitrary block shift.

It is nevertheless not the TFPT $\lambda$ field.  The source's common quarter
twist on eight complex edge channels has minimal charge $q=(1/4)^8$, hence
$h_q=1/4$, while the sourced $\lambda=(\omega_s,\omega_f)$ has
$\norm{\lambda}^2=2$ and $h_\lambda=1$.  Forcing weight one by an integer
dressing inside the common scalar phase class requires odd fermion charge.
The actual $A_3$ fundamental makes the missing datum explicit: it is a
four-channel quarter disorder plus one endpoint fermion, combined with five
half-twisted $D_5$ channels and a nontrivial cocycle/GSO identification.
Those species-resolved couplings and that GSO map are not present in the
audited QWZ path.  A sharp endpoint insertion also has measured $O(1)$ lattice
excitation energy; a smooth whole-circle dressing lowers this to $O(N^{-1})$
but loses sharp support.  The result is therefore a positive microscopic
deck-disorder construction and a narrow obstruction to identifying it with
$\Psi_\lambda$, not a no-go for a future species-resolved extension.  No RH
claim is made.
\end{keybox}

\paragraph{Frozen source and target.}

The construction below uses the hopping matrix literally implemented in
\path{verification/v988_psi_lambda_reduction.py}.  With
\[
 T_x=\frac{\sigma_x}{2i}-\frac{\sigma_z}{2},\qquad
 T_y=\frac{\sigma_y}{2i}-\frac{\sigma_z}{2},
\]
the one-particle QWZ Hamiltonian $h_r$ has mass $M=1$, open $y$ boundary,
and on the bond $(N-1,y)\to(0,y)$ the forward hopping is $i^rT_x$.
Every other longitudinal bond has $T_x$.  The executable v988 campaign
constant is $N_y=6$, but its builder accepts $N_y$ as an argument.  This
artifact calls that builder explicitly at the charged-limit manuscript value
$N_y=8$.  The manuscript writes
\[
 p_n^{(r)}=\frac{2\pi}{N}\left(n+\frac r4\right).
\tag{1.1}
\]
That is its opposite-sign Bloch convention.  With the checker convention
$\psi_x=e^{ikx}$ and the executable forward seam amplitude $i^r$, the
boundary eigencondition is
\[
 e^{i(r\pi/2+kN)}=1,\qquad
 k_n^{(r)}=\frac{2\pi n-r\pi/2}{N}.
\tag{1.2}
\]
The checker reconstructs the complete position-space matrix from these Bloch
blocks and compares it entrywise with the imported v988 builder.  Using the
plus sign in this checker convention fails that matrix test.
It also states that sixteen Majorana copies give eight complex edge bosons.

The target is not specified merely by the deck phase.  The repository fixes
\[
 \omega_s=\left(\tfrac12,\tfrac12,\tfrac12,\tfrac12,\tfrac12\right),
 \qquad
 \omega_f=\left(\tfrac34,-\tfrac14,-\tfrac14,-\tfrac14\right),
 \qquad
 \lambda=(\omega_s,\omega_f),
\tag{1.3}
\]
so
\[
 \norm{\omega_s}^2=\frac54,\quad
 \norm{\omega_f}^2=\frac34,\quad
 \norm{\lambda}^2=2,\quad h_\lambda=1.
\tag{1.4}
\]
These exact claims occur in two frozen sources:
\begin{itemize}[nosep,leftmargin=2em]
\item \path{verification/v983_simple_current_generator.py};
\item \path{articles/2026-08-30/mmst_charged_scaling_limit_en.tex}.
\end{itemize}
The latter still lists a support-preserving microscopic realization and an
independent FE-GEN or gauge/code-projected algebra theorem as open.  We do not
silently replace (1.3) by an arbitrary representative of a scalar $\Z_4$
phase.

\paragraph{A common microscopic Hilbert tied to the hopping.}

Let $\mathfrak h_N=\C^{2NN_y}$ and let $\F_N=\F(\mathfrak h_N)$ be the
canonical number-conserving CAR Fock space.  Adjoin one four-level link
register $\mathcal K=\C^4$ at the chosen periodic seam.  In the basis
$\{|r\rangle:r\in\Z_4\}$ define
\[
 Z|r\rangle=i^r|r\rangle,\qquad X|r\rangle=|r+1\rangle,
 \qquad ZXZ^*=iX,\qquad X^4=Z^4=\id.
\tag{2.1}
\]
If $F=\sum_y |0,y\rangle T_x\langle N-1,y|$ is the forward seam hopping and
$h_{\rm cut}=h_0-F-F^*$, set
\[
 \widetilde h
 =\id_{\mathcal K}\otimes h_{\rm cut}+Z\otimes F+Z^*\otimes F^*
 =\bigoplus_{r=0}^3 h_r.
\tag{2.2}
\]
Thus the register is not a decorative sector label: its clock occurs in the
actual QWZ hopping term.  On $\mathcal K\otimes\F_N$ take
\[
 \widetilde H=\sum_{r=0}^3 |r\rangle\langle r|
 \otimes d\Gamma(h_r).
\tag{2.3}
\]
The four sectors now live in one canonical Hilbert space.  This is an
\emph{auxiliary $\Z_4$ flux link}; no claim is made that the present TFPT
code derives it.

\paragraph{The sharp line-ended operator.}

Choose an arc $A=\{0,1,\ldots,\ell\}$ with $0\le\ell<N-1$.  Let $u_A$ be
the one-particle onsite gauge
\[
 u_A|x,y,a\rangle=
 \begin{cases}
 i|x,y,a\rangle,&x\in A,\\
 |x,y,a\rangle,&x\notin A,
 \end{cases}
 \qquad
 \Gamma(u_A)=\exp\!\left(\frac{\pi i}{2}N_A\right).
\tag{2.4}
\]
Define
\[
 D_A=X\otimes\Gamma(u_A).
\tag{2.5}
\]

\begin{theorem}[Microscopic QWZ disorder transporter]
For every $r\in\Z_4$,
\[
 P_{r+1}D_AP_r
 =|r+1\rangle\langle r|\otimes\Gamma(u_A)\ne0,
 \qquad D_A^4=\id.
\tag{2.6}
\]
Writing $\Pi=(-1)^{N_{\rm F}}$ for canonical total fermion parity,
\[
 [D_A,\id\otimes\Pi]=0,\qquad
 (Z\otimes\id)D_A(Z^*\otimes\id)=iD_A.
\tag{2.7}
\]
Moreover
\[
 [\widetilde H,D_A]P_r
 =|r+1\rangle\langle r|\otimes
 d\Gamma\!\left(h_{r+1}-u_Ah_ru_A^*\right)\Gamma(u_A),
\tag{2.8}
\]
and $h_{r+1}-u_Ah_ru_A^*$ is supported only on the single longitudinal
bond $(\ell,y)\leftrightarrow(\ell+1,y)$, for all $y=0,\ldots,7$.
Its one-particle rank is $2N_y=16$ and its norm is $\sqrt2$.
\end{theorem}

\begin{proof}
The corner and charge equations follow from (2.1).  Since $u_A^4=\id$ and
$X$ acts on a different tensor factor, $D_A^4=\id$.  A number rotation is an
even CAR automorphism, proving the parity commutator.

The $(r+1,r)$ one-particle commutator block is
$h_{r+1}u_A-u_Ah_r=(h_{r+1}-u_Ah_ru_A^*)u_A$; second quantization gives
(2.8).  Conjugation by $u_A$ multiplies a forward hopping $x\to x+1$ by
$u_{x+1}u_x^*$.  At the original seam this factor is $i$, changing $i^r$
to $i^{r+1}$.  All bonds internal to $A$ or its complement are unchanged.
Only the other boundary of $A$ has factor $-i$.  Hence the original seam
defect is transported to the far endpoint.  Substitution of $T_x$ gives
rank $16$ and norm $\sqrt2$; the companion probe checks the full matrices
for all four sectors at $N_y=8$.
\end{proof}

This is the promised positive result.  A matter-only onsite gauge could not
have done it: the product of link phases telescopes,
\[
 \prod_x a_x\frac{g_{x+1}}{g_x}=\prod_xa_x,
\tag{2.9}
\]
so it preserves the Wilson holonomy.  The new register is exactly the extra
degree of freedom that changes it.  Unlike the Round-6 arbitrary $J_r$, the
operator (2.5) is tied to (2.2), has a specified CAR action, and obeys the
endpoint commutator identity (2.8).

Equation (2.7) proves compatibility with canonical fermion parity and with
any GSO projector that is a function of that parity.  It does \emph{not}
prove compatibility with the repository's character-level $O+S$ selector:
the audited source supplies no microscopic projector identifying that
selector with $\Pi$.  This boundary is essential.

\paragraph{Why the constructed operator is not \texorpdfstring{$\Psi_\lambda$.}{Psi-lambda}}

For a normalized complex chiral fermion, a twist by $e^{2\pi iq}$ has
minimal bosonized weight $q^2/2$, with $q\in(-1/2,1/2]$.  The actual common
phase $i$ therefore fixes, for the eight complex channels in the source,
\[
 q_{\rm deck}=\left(\tfrac14,\ldots,\tfrac14\right),\qquad
 \norm{q_{\rm deck}}^2=\frac12,\qquad h_{\rm deck}=\frac14.
\tag{3.1}
\]
It has order four but nontrivial spin $e^{2\pi ih}=i$.  The target (1.3)
instead has integer spin and weight one.  Norm is invariant under an
orthonormal change of current basis, so this is not repaired by relabeling
the eight bosons.

The phase by itself does not determine an excited representative.  Indeed
\[
 q'=\left(\tfrac54,\tfrac14,\ldots,\tfrac14\right)
\tag{3.2}
\]
has the same scalar phase and $\norm{q'}^2=2$.  But $q'-q_{\rm deck}$ is one
unit of fermion charge.  More generally, if $q_{\rm deck}+n$ has norm two,
then
\[
 2\left(\norm{q_{\rm deck}+n}^2-\norm{q_{\rm deck}}^2\right)
 =2\sum_jn_j^2+\sum_jn_j=3.
\tag{3.3}
\]
Modulo two, $\sum_jn_j$ is odd.  Every such weight-one dressing therefore
changes canonical fermion parity.  Selecting (3.2) solely because it has the
desired number would be precisely an arbitrary charge-basis shift, not a
microscopic derivation of $\lambda$.

The $D_5\oplus A_3$ coordinates show the missing operator content more
constructively.  The $D_5$ spinor part is a half disorder on five complex
channels and contributes $5/8$.  For the $A_3$ part,
\[
 \left(\tfrac34,-\tfrac14,-\tfrac14,-\tfrac14\right)
 =\left(-\tfrac14,-\tfrac14,-\tfrac14,-\tfrac14\right)+(1,0,0,0).
\tag{3.4}
\]
Thus it is a four-channel quarter disorder of weight $1/8$ plus one endpoint
fermion, giving the traceless $A_3$ fundamental weight $3/8$.  Combined with
the five half twists, this gives $h_\lambda=5/8+3/8=1$.

Equation (3.4) is the next microscopic specification.  Indeed,
$e^{2\pi i\omega_s}=(-1)^5$ componentwise and
$e^{2\pi i\omega_f}=(-i)^4$ componentwise.  A $\lambda$ operator therefore
needs species-resolved seam clocks $Z^2$ on the five $D_5$ channels and
$Z^{-1}$ on the four $A_3$ fermions, the endpoint fermion, removal of the
diagonal $U(1)$ in the $A_3$ realization, and the cocycle/GSO rule that makes
the total field local and bosonic.  The frozen QWZ code has one common scalar
$i^r$ phase and none of these four identification maps.  The flux register
solves the common-Hilbert and locality problem, but not this internal-charge
problem.  This redundant fermionic realization has nine complex fermions
before the $A_3$ diagonal-$U(1)$ quotient; it is therefore not automatically
an operator on the source's eight-channel Fock space.  An explicit conformal
embedding/isometry is part of the missing microscopic datum.

\paragraph{Energy and support diagnostics.}

There is a finite-lattice tradeoff that prevents another overclaim.
Applying the sharp $D_A$ to the half-filled $r=0$ ground state and measuring
above the $r=1$ ground energy gives, for $N=16,24,32,48,64$,
\[
 4.774031,\quad4.730808,\quad4.708887,\quad4.686751,\quad4.675602.
\tag{4.1}
\]
The cost is $O(1)$ at fixed $N_y=8$, consistent with the sharp endpoint
commutator.  This bare operator is not a low-energy weight-one state.

One can distribute the string by taking
\[
 u_N(x)=\exp\!\left[\frac{\pi i}{2}\left(1-\frac{x}{N}\right)\right],
 \qquad D_N^{\rm sm}=X\otimes\Gamma(u_N).
\tag{4.2}
\]
Then the actual one-particle matrices give
\[
 N\norm{h_1u_N-u_Nh_0}\longrightarrow\frac\pi2,\qquad
 N\left(\langle D_N^{\rm sm}\Omega_0,H_1D_N^{\rm sm}\Omega_0\rangle
 -E_0(H_1)\right)\approx8.475.
\tag{4.3}
\]
This is measured $O(N^{-1})$ energy, but $u_N$ occupies the whole
circumference and the coefficient contains full-band polarization.  It is
not a proof of conformal weight.  Conversely, the sharp string has exact
endpoint support but the $O(1)$ UV cost (4.1).  Recovering both sharp local
algebras and the weight-one edge primary requires the missing edge projection,
renormalized endpoint field, species-resolved charge, and GSO cocycle.

As an independent finite diagnostic, momentum-block diagonalization gives
\[
 \frac{N}{2\pi}\bigl(E_0(h_1)-E_0(h_0)\bigr)\longrightarrow-\frac3{16}
\tag{4.4}
\]
for one two-edge QWZ copy.  This is a sector-vacuum/Casimir shift, not the
one-edge primary weight and not $h_\lambda$.  No fit or finite sample in this
section is promoted to an all-$N$ theorem.

\paragraph{Executable certificate and claim boundary.}

The repository certificate \veri{v1033\_\allowbreak charged\_\allowbreak disorder.py}
audits the frozen sources, proves the finite Weyl/parity/charge identities
with rational and symbolic arithmetic, imports the executable v988 QWZ
builder and compares it entrywise, verifies the endpoint commutator on the
actual $M=1,N_y=8$ hops, and reports finite energy diagnostics.

\begin{center}
\fbox{\parbox{0.91\linewidth}{
\textbf{Established:} after adjoining one explicitly auxiliary local
$\Z_4$ flux link, $D_A=X\Gamma(u_A)$ is a CAR-even, order-four, grade-one
microscopic QWZ boundary-condition changer with an exact line-endpoint
commutator.\\[2pt]
\textbf{Not established:} that the present TFPT/GSO construction contains
this register; that its $O+S$ selector is canonical parity; or that the
common quarter disorder equals the $D_5\oplus A_3$ field $\Psi_\lambda$.
The latter identification is false for the minimal common twist because
$h_{\rm deck}=1/4\ne1=h_\lambda$.}}
\end{center}
\endgroup
